\section{A uniform local finite-Gauss-fibre sector}
\begin{proposition}\label{uniformfibres}
At any full point $x$, the finite-flat representative used in
\eqref{atom} supplies a compact set $\Theta$ containing $x$, formed
by all points of finitely many fibres above a small compact family of
weighted Gauss points. There is a uniformly bounded family of global
caps $w_r\in\C$ with uniformly bounded normalized energies and
\[
 \MA(w_r)\ge\epsilon^2\delta_y
 \quad\text{at EVERY }y\text{ of its corresponding finite Gauss fibre},
 \qquad\epsilon>0\text{ independent of }r,y.
\]
Every normalized full-dual maximizing net converges uniformly on
$\Theta$ to its finite continuous full-point limit. For arbitrarily
small such parameter families all their fibres concentrate near $x$.
This supplies no affinoid-neighborhood basis with all Shilov boundaries
in $\Theta$.
\end{proposition}
\begin{proof}
Use the two fixed Laurent polynomials $P_1,P_2$ chosen at $x$ and a
compact finite-flat representative $f:V\to U$ of positive degree $D$,
with singleton fibre at $g_0=f(x)$, a weighted Gauss point. The
representative domains are analytic neighborhoods, so $x\in\operatorname{Int}_X V$
and $g_0\in\operatorname{Int}U$. Shrink the parameter rectangle $J$
around its log radii until its continuous Gauss section $\sigma(J)$
lies in $\operatorname{Int}U$ and
$\Theta=f^{-1}(\sigma(J))\cap V$ lies in the ordinary interior of
$V$ and in a chosen small base-coordinate ball around $x$.
Indeed, for any ordinary neighborhood $O$ of $x$, $V\setminus O$
is compact and its image is compact, disjoint from $g_0$ by the singleton
fibre. A target neighborhood of $g_0$ avoids this image; continuity of
$\sigma$ gives the required $J$. This also proves concentration for
arbitrary $O$. Compactness of $V$ and of $\sigma(J)$ gives compact
$\Theta$. All these points have full graded dimension, since finite
complete residue-field extensions of a two-dimensional weighted Gauss
field preserve that dimension. No varying section is asserted analytic.

On a larger injective base tropical ball of radius $R$, set
$\rho_0=|\tau-\tau(x)|^2/2$ and
$\psi_r=\log\max(1,|P_1|/r_1,|P_2|/r_2)$.
The compact larger polyannulus bounds $|P_i|$, and the compact parameter
rectangle bounds each $r_i$ away from zero, so $0\le\psi_r\le M_0$
uniformly. Choose one $0<\delta<R^2/4$ and one $\epsilon>0$ with
$\delta+\epsilon M_0$ smaller than the minimum of $\rho_0$ on a
fixed outer collar. The caps
$w_r=\max(0,\delta+\epsilon\psi_r-\rho_0)$ glue by zero to
actual continuous weights at all stages, including at polynomial zeros,
by the ramp/maximum construction already proved. Their upper bound
$M=\delta+\epsilon M_0$ is uniform. Shrink $J$ by the preceding
properness argument so $\Theta$ lies where $\rho_0<\delta/2$;
there $\theta_{w_r}=\epsilon\DD\psi_r$.

At every point $y$ of the finite Gauss fibre, CLD18.6.2 identifies the
top log-max product with its actual fibre integration current. Local
finite-flat calibration14.5.2--3 makes its point coefficient the positive
integer local degree $d_y\ge1$. This is the same degree computation as
\eqref{atom}, applied at all points of this finite fibre; the sum of
local degrees is $D$. Thus the cap coefficient is
$\epsilon^2d_y\ge\epsilon^2$. All domains here are restricted to
ordinary interiors for the boundaryless current formula. There is no
independent field-extension factor in addition to this local degree.
Normalize $\widehat w_r=w_r-\sup_Xw_r$.
Then $-M\le\widehat w_r\le0$ and monotonicity gives
$E(\widehat w_r)\ge-2M$ uniformly.

For $a,b\in\C_C$, put $D_0=\max(C,2M,0)$.
The fixed-test estimate of Theorem \ref{nonorder}, its proof, and
\eqref{mom} give uniformly in $r$
\[
 \int|a-b|\MA(w_r)
 \le\Omega_C(I)+2\sqrt{6D_0}\sqrt{\Psi_{D_0}(I)}.
\]
The atom bound divides this by $\epsilon^2$ to control
$|a(y)-b(y)|$ at every $y\in\Theta$, independently of $r,y$.
Equation \eqref{deficit} tends to zero for maximizing tails, proving
uniform Cauchyness on $\Theta$. Each restriction is continuous;
completeness of $C(\Theta,\R)$ gives a finite continuous limit equal
to the full-point tuple. Its continuity and the properness concentration
prove convergence of its values along these fibres towards its value at $x$.
\end{proof}

There is a genuine unresolved next step. A neighborhood of a target
Gauss point can exclude cancellation points of the same coordinate
radii. Therefore a shrinking coordinate polyannulus need not be
contained in the target neighborhood $U$ above. Pulling back such a
polyannulus is not a justified affinoid-neighborhood construction for
this local representative. The assertion that all required nearby
Shilov boundaries lie in $\Theta$ needs another argument; properness
and the uniform cap estimate alone do not prove it. In particular this
proposition is not used to assert $R(x)=\ell(x)$ or positivity of $R$.

